% Lösungen Einheit 1 von 5 – Winkel messen und zeichnen (zusammenbau v0.1)
\einheitenkopf{Einheit 1 von 5 · Winkel messen und zeichnen}

\erg{8}{a) die kleinere Zahl}
%% TODO Lösung der Erklärzeile 9g) fehlt in der Bank
\erg{9}{a) spitz \quad b) $\alpha = 37^\circ$ \quad c) $\alpha = 54^\circ$ \quad d) $\alpha = 18^\circ$ \quad e) $\alpha = 83^\circ$ \quad f) $\alpha = 66^\circ$ \quad g) \ldots \quad h) $\alpha = 118^\circ$ \quad i) Nullpunkt auf S, Nulllinie auf dem Strahl, bei $55^\circ$ markieren, zweiten Schenkel von S aus durch die Markierung zeichnen \quad j) Rest $145^\circ$; $\alpha = 360^\circ - 145^\circ = 215^\circ$ \quad k) $\angle BAC$ (oder $\angle CAB$), Scheitel A \quad l) $34^\circ + 47^\circ = 81^\circ$ \quad m) $96^\circ - 38^\circ = 58^\circ$ \quad n) 1,2 km = 1\,200 m; 1\,200 m − 450 m = 750 m \quad o) $\gamma = 44^\circ - 41^\circ = 3^\circ$}
\erg{10}{a) ein Viertel}
\erg{11}{a) Fehler: falsche Zahlenreihe – α ist spitz; richtig: $\alpha = 180^\circ - 144^\circ = 36^\circ$ \quad b) Die Grenze ist der rechte Winkel: $91^\circ$ ist größer als ein rechter, $89^\circ$ kleiner. \quad c) $90^\circ - 54^\circ = 36^\circ$}
